\documentclass[11pt,a4paper]{article}

\usepackage[a4paper,margin=14mm,top=15mm,bottom=16mm]{geometry}
\usepackage{fontspec}
\usepackage{polyglossia}
\setmainlanguage{russian}
\setotherlanguage{english}
\setmainfont{Arial Unicode MS}
\newfontfamily\cyrillicfont{Arial Unicode MS}
\newfontfamily\cyrillicfontsf{Arial Unicode MS}
\newfontfamily\headingfont{Georgia}
\newfontfamily\sansfont{Arial Unicode MS}
\usepackage[table]{xcolor}
\usepackage{tikz}
\usetikzlibrary{calc,positioning}
\usepackage[skins,breakable]{tcolorbox}
\usepackage{tabularx}
\usepackage{booktabs}
\usepackage{array}
\usepackage{enumitem}
\usepackage{multicol}
\usepackage{parskip}
\usepackage{fancyhdr}
\usepackage{hyperref}
\usepackage{graphicx}
\usepackage{etoolbox}

\hypersetup{hidelinks}

\definecolor{paperbg}{HTML}{F5EFE7}
\definecolor{paperwhite}{HTML}{FFFDFB}
\definecolor{ink}{HTML}{1D2332}
\definecolor{muted}{HTML}{60697A}
\definecolor{navy}{HTML}{20283D}
\definecolor{navysoft}{HTML}{39435F}
\definecolor{accent}{HTML}{C1874F}
\definecolor{accentsoft}{HTML}{F2E4D6}
\definecolor{emerald}{HTML}{1F6B62}
\definecolor{emeraldsoft}{HTML}{E2EFEA}
\definecolor{rose}{HTML}{8E4A5D}
\definecolor{rosesoft}{HTML}{F4E6EB}
\definecolor{line}{HTML}{D7D2CB}

\pagecolor{paperbg}
\color{ink}

\setlength{\parindent}{0pt}
\setlength{\parskip}{4pt}
\setlength{\emergencystretch}{2em}
\renewcommand{\arraystretch}{1.28}

\pagestyle{fancy}
\fancyhf{}
\fancyfoot[R]{\color{muted}\sffamily\small \thepage}
\renewcommand{\headrulewidth}{0pt}

\newcolumntype{Y}{>{\raggedright\arraybackslash}X}

\tcbset{
  frame hidden,
  sharp corners,
  boxrule=0pt,
  enhanced,
  breakable
}

\newtcolorbox{glassbox}[1][]{
  colback=paperwhite,
  borderline={0.5pt}{0pt}{line},
  arc=7mm,
  left=5mm,
  right=5mm,
  top=4mm,
  bottom=4mm,
  drop shadow={black!12!paperbg},
  #1
}

\newtcolorbox{softbox}[1][]{
  colback=white,
  borderline={0.5pt}{0pt}{line},
  arc=5mm,
  left=4mm,
  right=4mm,
  top=3mm,
  bottom=3mm,
  #1
}

\newtcolorbox{navybox}[1][]{
  colback=navy,
  coltext=white,
  arc=7mm,
  left=5mm,
  right=5mm,
  top=4mm,
  bottom=4mm,
  drop shadow={black!22!paperbg},
  #1
}

\newcommand{\eyebrow}[1]{{\sansfont\fontsize{8.5}{10}\selectfont\color{muted}\MakeUppercase{#1}}}
\newcommand{\titlefont}[1]{{\headingfont #1}}
\newcommand{\pill}[1]{%
  \tikz[baseline=(n.base)] \node[rounded corners=9pt, fill=white, fill opacity=0.08, text opacity=1,
  draw=white!18, inner xsep=9pt, inner ysep=5.5pt] (n)
  {\color{white}\sansfont\fontsize{8.5}{9}\selectfont\MakeUppercase{#1}};%
}
\newcommand{\kicker}[2]{%
  \begin{softbox}[colback=#1!10]
    {\sansfont\bfseries\color{#1} #2}
  \end{softbox}
}
\newcommand{\stephead}[2]{%
  \begin{tikzpicture}[baseline=(title.base)]
    \node[rounded corners=4mm, fill=navy, minimum width=12mm, minimum height=12mm, text=white, font=\headingfont\bfseries\large] (badge) {#1};
    \node[right=4mm of badge, anchor=west] (title) {\begin{minipage}{0.8\linewidth}{\headingfont\fontsize{18}{20}\selectfont #2}\end{minipage}};
  \end{tikzpicture}
}
\newcommand{\minititle}[1]{{\sansfont\bfseries\fontsize{9.5}{11}\selectfont\color{rose}\MakeUppercase{#1}}}

\begin{document}
\thispagestyle{empty}

\begin{titlepage}
\pagecolor{navy}
\begin{tikzpicture}[remember picture,overlay]
  \fill[navy] (current page.south west) rectangle (current page.north east);
  \fill[accent, opacity=0.16] ($(current page.north west)+(2.2,-2.0)$) circle (3.2cm);
  \fill[emerald, opacity=0.14] ($(current page.north east)+(-2.1,-2.6)$) circle (2.4cm);
  \fill[accent, opacity=0.11] ($(current page.south west)+(1.8,1.6)$) circle (2.7cm);
  \draw[white, opacity=0.10, line width=0.5pt] ($(current page.north east)+(-2.1,-2.6)$) circle (3.8cm);
  \draw[white, opacity=0.07, line width=10pt] ($(current page.north east)+(-2.1,-2.6)$) circle (3.0cm);
  \draw[white, opacity=0.05, line width=18pt] ($(current page.north east)+(-2.1,-2.6)$) circle (2.1cm);
\end{tikzpicture}

\vspace*{8mm}
\noindent\pill{Confidential \hspace{2mm} Переговорный playbook}
\hfill
\pill{@monicastyle.mc \hspace{2mm} April 2026}

\vspace{24mm}

\begin{minipage}[t]{0.62\linewidth}
{\color{white!68}\sansfont\fontsize{9}{11}\selectfont FBI Negotiation + NLP Frame\par}
\vspace{6mm}
{\headingfont\color{white}\fontsize{28}{30}\selectfont Playbook защиты команды @monicastyle.mc\par}
\vspace{8mm}
{\color{white!82}\fontsize{13.2}{18}\selectfont
Документ для спокойных и уверенных переговоров с клиентом, который хочет сократить бюджет.
Цель - не спорить, а мягко вернуть разговор к ценности системы, ролей и 8-недельного спринта.\par}

\vspace{14mm}

\begin{tcolorbox}[colback=white, colframe=white, opacityback=0.08, opacityframe=0.12, arc=6mm, left=4mm, right=4mm, top=3mm, bottom=3mm]
{\sansfont\bfseries\color{white!92} Ситуация}\par
{\color{white!76}\fontsize{10.8}{14}\selectfont Клиент хочет оставить только монтажёра, SMM и PM, потому что пока не видит достаточный ROI на ранней стадии аккаунта.}
\end{tcolorbox}

\vspace{4mm}

\begin{tcolorbox}[colback=white, colframe=white, opacityback=0.08, opacityframe=0.12, arc=6mm, left=4mm, right=4mm, top=3mm, bottom=3mm]
{\sansfont\bfseries\color{white!92} Наша задача}\par
{\color{white!76}\fontsize{10.8}{14}\selectfont Сохранить команду через структуру Chris Voss, маркировку эмоций клиента и NLP-рефрейм ролей как функций бренда.}
\end{tcolorbox}

\vspace{4mm}

\begin{tcolorbox}[colback=white, colframe=white, opacityback=0.08, opacityframe=0.12, arc=6mm, left=4mm, right=4mm, top=3mm, bottom=3mm]
{\sansfont\bfseries\color{white!92} Тон переговоров}\par
{\color{white!76}\fontsize{10.8}{14}\selectfont Late-Night FM DJ: медленно, тёпло, уверенно. Никакого давления, только ясность, факты и чувство контроля у клиента.}
\end{tcolorbox}
\end{minipage}
\hfill
\begin{minipage}[t]{0.32\linewidth}
\vspace{18mm}
\begin{tcolorbox}[colback=white, colframe=white, opacityback=0.09, opacityframe=0.14, arc=9mm, left=6mm, right=6mm, top=8mm, bottom=8mm, height=0.56\textheight]
{\color{white}\fontsize{36}{36}\selectfont 01\par}
\vspace{18mm}
{\color{white!70}\sansfont\fontsize{8.5}{10}\selectfont Ключевой тезис\par}
\vspace{3mm}
{\headingfont\color{white}\fontsize{18}{20}\selectfont Не 7 человек. 4 функции бренда.\par}
\vspace{7mm}
{\color{white!80}\fontsize{11.2}{15}\selectfont
Убрать людей легко. Восстановить систему, визуальное превосходство, reach и темп роста - в разы дороже.\par}
\vfill
{\color{white!68}\sansfont\fontsize{8.5}{11}\selectfont Основа документа\par}
{\color{white!78}\fontsize{10.2}{14}\selectfont Chris Voss, NLP framework, Reddit / X research, маркетинговые наблюдения 2026.}
\end{tcolorbox}
\end{minipage}
\end{titlepage}

\clearpage
\pagecolor{paperbg}
\color{ink}
\setcounter{page}{2}

\begin{navybox}
\eyebrow{Стратегическая рамка}\par
\vspace{1.5mm}
{\fontsize{13.4}{17}\selectfont \textbf{Золотое правило Voss:} люди давят не потому, что правы, а потому что чувствуют, что их не слышат. Убери это чувство - и давление исчезает.}
\end{navybox}

\vspace{4mm}

\begin{glassbox}
{\headingfont\fontsize{22}{24}\selectfont Executive Summary}\hfill
{\color{muted}\sansfont\fontsize{9.5}{12}\selectfont Сначала снимаем защиту, потом даём клиенту ощущение контроля, потом продаём не людей, а систему.}

\vspace{4mm}

\begin{tabularx}{\linewidth}{>{\raggedright\arraybackslash}p{0.22\linewidth} >{\raggedright\arraybackslash}p{0.22\linewidth} >{\raggedright\arraybackslash}p{0.22\linewidth} >{\raggedright\arraybackslash}p{0.22\linewidth}}
\rowcolor{paperbg}
\textbf{279} & \textbf{8 недель} & \textbf{4 функции} & \textbf{2x} \\
\fontsize{9.6}{12}\selectfont Подписчиков сейчас. Клиент видит это как ``слишком мало для такой команды''. &
\fontsize{9.6}{12}\selectfont Реалистичный горизонт, чтобы честно оценить рост и не ломать систему преждевременно. &
\fontsize{9.6}{12}\selectfont Визуал, рост, credibility, системность. Именно их нужно защищать. &
\fontsize{9.6}{12}\selectfont Ориентир стоимости перезапуска после паузы по сравнению с удержанием темпа сейчас. \\
\end{tabularx}

\vspace{5mm}

\begin{tabularx}{\linewidth}{Y Y}
\rowcolor{paperbg}
\textbf{1. Accusation Audit} & \textbf{2. Labeling} \\
\fontsize{9.8}{12}\selectfont Озвучиваем все страхи клиента первыми и гасим защитную реакцию. &
\fontsize{9.8}{12}\selectfont Называем чувства вслух, чтобы снизить напряжение и вернуть доверие. \\
\rowcolor{paperbg}
\textbf{3. That's Right} & \textbf{4. Reframing} \\
\fontsize{9.8}{12}\selectfont Точно суммируем позицию клиента и добиваемся ``Вот именно''. &
\fontsize{9.8}{12}\selectfont Меняем угол зрения: не ``много людей'', а ``функции бренда''. \\
\rowcolor{paperbg}
\textbf{5. Data Layer} & \textbf{6. Team Defense} \\
\fontsize{9.8}{12}\selectfont Подкрепляем разговор фактами про алгоритм, CAC, mindshare и скорость роста. &
\fontsize{9.8}{12}\selectfont Каждый человек получает конкретную функцию и конкретную цену потери. \\
\rowcolor{paperbg}
\textbf{7. Future Pacing} & \textbf{8. Agreement Frame} \\
\fontsize{9.8}{12}\selectfont Рисуем два будущих: с системой и с урезанной командой. &
\fontsize{9.8}{12}\selectfont Продаём 8-недельный спринт с KPI вместо прямого отказа на сокращение. \\
\end{tabularx}
\end{glassbox}

\vspace{4mm}

\begin{multicols}{2}
\begin{glassbox}
{\headingfont\fontsize{18}{20}\selectfont Что клиент думает}\par
\vspace{2mm}
\begin{softbox}[colback=accentsoft]
\begin{itemize}[leftmargin=5mm,itemsep=2.5mm,topsep=1mm]
  \item ``Слишком дорого для аккаунта с 279 подписчиками.''
  \item ``Я могу справиться меньшей командой.''
  \item ``Не вижу ROI от каждого человека.''
  \item ``Пока нет результата - зачем платить за большую команду?''
\end{itemize}
\end{softbox}
\end{glassbox}

\columnbreak

\begin{glassbox}
{\headingfont\fontsize{18}{20}\selectfont Что за этим стоит}\par
\vspace{2mm}
\begin{softbox}[colback=emeraldsoft]
\minititle{Страх}\par
Деньги уходят быстрее, чем виден рост. Это про тревогу, а не про враждебность.
\end{softbox}
\vspace{2.5mm}
\begin{softbox}[colback=emeraldsoft]
\minititle{Контроль}\par
Клиент хочет понимать, за что платит, и ощущать измеримость вложений.
\end{softbox}
\vspace{2.5mm}
\begin{softbox}[colback=emeraldsoft]
\minititle{Ранняя стадия}\par
Проект ещё не доказал себя. Значит, нужно продавать не прошлый успех, а систему будущего.
\end{softbox}
\vspace{2.5mm}
\begin{softbox}[colback=emeraldsoft]
\minititle{Внешнее давление}\par
Могут влиять партнёр, семья или финансы. Это надо уважить, а не спорить с этим.
\end{softbox}
\end{glassbox}
\end{multicols}

\clearpage

\begin{glassbox}
\stephead{1}{Accusation Audit}

\vspace{3mm}
\begin{softbox}[colback=accentsoft]
\textbf{Цель:} перечислить вслух все возможные обвинения до того, как клиент их сформулирует. Это снимает напряжение и переводит разговор из конфликта в сотрудничество.
\end{softbox}

\vspace{3mm}
\begin{navybox}
\eyebrow{Скрипт Яны}\par
\vspace{2mm}
``Прежде чем мы начнём обсуждать дальнейшие шаги, мне важно сказать кое-что.
Я уверена, ты думаешь, что команда слишком большая для аккаунта с 279 подписчиками.
Возможно, тебе кажется, что мы перестарались и что всё это можно сделать проще и дешевле.
Вполне вероятно, ты чувствуешь, что платишь за много людей, а результат пока не на том уровне, на котором хочется его видеть.
И может быть, ты даже думаешь, что мы просто оправдываем существование каждой позиции, чтобы удержать контракт.
Я хочу, чтобы ты знала - мы это понимаем. И я не буду с тобой спорить. Вместо этого я хочу показать тебе, что именно делает каждый человек и что случится с проектом, если его убрать.'' 
\end{navybox}

\vspace{3mm}
\begin{tabularx}{\linewidth}{Y Y}
\rowcolor{paperbg}
\textbf{Что произойдёт} & \textbf{Что важно не делать} \\
\fontsize{10}{13}\selectfont Клиент почти автоматически начнёт смягчать свою позицию: ``Ну, я не прямо так считаю...'' Именно это и нужно. Защита уже ослаблена. &
\fontsize{10}{13}\selectfont Не оправдываться, не спорить и не торопиться с цифрами. Сначала нужно дать клиенту почувствовать, что вы уже услышали её страхи. \\
\end{tabularx}
\end{glassbox}

\vspace{4mm}

\begin{glassbox}
\stephead{2}{Labeling}

\vspace{3mm}
\begin{softbox}[colback=accentsoft]
\textbf{Цель:} вслух маркировать переживание клиента словами ``Похоже...'' или ``Кажется...'', чтобы она почувствовала: её состояние поняли без нападения.
\end{softbox}

\vspace{3mm}
\rowcolors{2}{paperbg}{paperwhite}
\begin{tabularx}{\linewidth}{>{\raggedright\arraybackslash}p{0.31\linewidth} Y}
\toprule
\rowcolor{paperbg}
\textbf{Когда клиент говорит} & \textbf{Лейбл в ответ} \\
Мне кажется, это слишком дорого. & Похоже, ты чувствуешь, что инвестиция не пропорциональна тому, где аккаунт находится сейчас. \\
Я хочу оставить только троих. & Кажется, тебе важно чувствовать, что каждый человек реально нужен, а не просто присутствует в команде. \\
279 подписчиков - это мало. & Звучит так, будто ты переживаешь, что рост идёт медленнее, чем ты ожидала. \\
Мне посоветовали сократить. & Похоже, на тебя есть внешнее давление, и тебе важно принять максимально взвешенное решение. \\
\bottomrule
\end{tabularx}
\end{glassbox}

\vspace{4mm}

\begin{glassbox}
\stephead{3}{That's Right}

\vspace{3mm}
\begin{navybox}
\eyebrow{Скрипт перехода}\par
\vspace{2mm}
``Правильно ли я понимаю? Ты видишь аккаунт с 279 подписчиками и при этом команду из 7 человек.
Ты чувствуешь, что соотношение не сходится. Ты хочешь видеть, что каждый рубль работает на рост, а не просто на процесс.
И при этом тебе важно не потерять визуальное качество, которое уже есть, потому что именно оно отличает Monica от остальных в этой нише.'' 
\end{navybox}

\vspace{3mm}
\begin{tabularx}{\linewidth}{Y Y}
\rowcolor{paperbg}
\textbf{Ожидаемый ответ} & \textbf{Почему это ключевой момент} \\
\fontsize{10}{13}\selectfont ``Вот именно'' или ``Да, именно так''. После этого окно к предложению открывается. &
\fontsize{10}{13}\selectfont До ``That's Right'' клиент отстаивает позицию. После ``That's Right'' клиент готов рассматривать путь вперёд. \\
\end{tabularx}
\end{glassbox}

\clearpage

\begin{glassbox}
\stephead{4}{Reframing}

\vspace{3mm}
\begin{softbox}[colback=accentsoft]
\textbf{Главный рефрейм:} не ``вы платите за 7 человек'', а ``вы удерживаете 4 функции бренда, без которых система не выдержит стадию роста''.
\end{softbox}

\vspace{3mm}
\rowcolors{2}{paperbg}{paperwhite}
\begin{tabularx}{\linewidth}{>{\raggedright\arraybackslash}p{0.31\linewidth} Y}
\toprule
\rowcolor{paperbg}
\textbf{Как клиент видит} & \textbf{Как мы переформатируем} \\
7 человек - дорого. & 4 функции бренда, без которых система разваливается. \\
279 подписчиков - мало. & 279 - это ядро. Мы на стартовой площадке ракеты, а не на финише. \\
Можно обойтись тремя. & Трое умеют поддерживать. Семеро строят бренд, который потом сможет зарабатывать. \\
Пока нет ROI. & Видимый ROI у этой системы проявляется через 8 недель. Ей нужно время, чтобы доказать себя. \\
\bottomrule
\end{tabularx}

\vspace{3mm}
\begin{navybox}
\eyebrow{Ключевая метафора}\par
\vspace{2mm}
``Представь, что Monica - это стартап. Сейчас ты на pre-seed стадии.
Все стартапы на этом этапе инвестируют больше, чем зарабатывают.
Но те, кто сокращают команду именно на стадии роста, потом теряют в два раза больше,
когда приходится нанимать заново и разгонять алгоритм практически с нуля.''
\end{navybox}
\end{glassbox}

\vspace{4mm}

\begin{glassbox}
\stephead{5}{Факты и данные}

\vspace{3mm}
\rowcolors{2}{paperbg}{paperwhite}
\begin{tabularx}{\linewidth}{>{\raggedright\arraybackslash}p{0.30\linewidth} >{\raggedright\arraybackslash}p{0.44\linewidth} >{\raggedright\arraybackslash}p{0.18\linewidth}}
\toprule
\rowcolor{paperbg}
\textbf{Факт} & \textbf{Зачем он нужен в разговоре} & \textbf{Источник} \\
Стоимость перезапуска может быть в 2x выше, чем стоимость поддержания системы. & Показывает, что сокращение сегодня может создать более дорогую проблему завтра. & The Drum, MarketingFuel UK \\
Instagram быстро ``забывает'' аккаунт после 2-3 недель непоследовательности. & Объясняет, почему ежедневный ритм и непрерывность нельзя недооценивать. & SpurNow, InfluenceFlow 2026 \\
CAC обычно растёт, когда маркетинговую систему режут слишком рано. & Показывает, что каждый новый подписчик потом может обходиться дороже. & CMSwire \\
Engagement velocity в первые 30-60 минут - ключевой сигнал алгоритму. & Поддерживает важность тайминга, сторис, комментариев и системного SMM. & Storybox HQ \\
Бренды, которые замолкают, могут терять до 50\% mindshare за 4 недели. & Добавляет вес аргументу о цене паузы и потере внимания аудитории. & One2Create UK \\
Нано-инфлюенсеры с ER 3.5-8\% - один из самых ценных сегментов для брендов в 2026. & Рамка ``279 подписчиков'' меняется на рамку ``высокий актив ранней стадии''. & Amplify Influencer \\
\bottomrule
\end{tabularx}

\vspace{2mm}
{\color{muted}\sansfont\fontsize{8.8}{11}\selectfont Данные в документе используются как разговорный слой поддержки рефрейма. Их задача - не ``забомбить цифрами'', а придать уверенность ключевому тезису: раннюю систему нельзя резать по метрике одной недели.}
\end{glassbox}

\clearpage

\begin{glassbox}
{\headingfont\fontsize{22}{24}\selectfont Защита команды - Часть I}\hfill
{\color{muted}\sansfont\fontsize{9.5}{12}\selectfont Каждый человек = отдельная функция системы. Убирая человека, клиент убирает не роль, а рычаг роста.}
\end{glassbox}

\vspace{3mm}

\begin{glassbox}
{\headingfont\fontsize{18}{20}\selectfont Ира}\hfill
{\sansfont\bfseries\color{emerald} Визуальный язык}

{\color{muted} Визуал, криейтор, идеи, креатив}

\vspace{2mm}
\begin{softbox}[colback=paperbg]
\minititle{Функция в системе}\par
Создаёт визуальный язык бренда Monica и удерживает ощущение quality overshoot - когда маленький аккаунт выглядит как бренд на 50K.
\end{softbox}
\vspace{2mm}
\begin{softbox}[colback=paperbg]
\minititle{Bullet Point для клиента}\par
``Ира делает так, что Monica на 279 подписчиках выглядит как аккаунт на 50K. Наш AI-аудит показал, что визуальное качество контента уже на уровне микро-инфлюенсера. Это не случайность. Если убрать Иру, контент станет обычным - и бренд потеряет своё главное конкурентное преимущество.''
\end{softbox}
\vspace{2mm}
\begin{tabularx}{\linewidth}{Y Y}
\rowcolor{paperbg}
\textbf{Калиброванный вопрос} & \textbf{NLP-привязка} \\
Как ты видишь сохранение этого визуального уровня без человека, который его создаёт? & Все 4 рубрики визуально завязаны на Иру. Без неё Monica Style Edit теряет эстетику ``styled for her story''. \\
\end{tabularx}
\end{glassbox}

\vspace{3mm}

\begin{glassbox}
{\headingfont\fontsize{18}{20}\selectfont Лера}\hfill
{\sansfont\bfseries\color{emerald} Алгоритм и рост}

{\color{muted} Head of SMM}

\vspace{2mm}
\begin{softbox}[colback=paperbg]
\minititle{Функция в системе}\par
Стратегический мозг роста: алгоритм, хештеги, timing, caption SEO, системный анализ сигналов reach.
\end{softbox}
\vspace{2mm}
\begin{softbox}[colback=paperbg]
\minititle{Bullet Point для клиента}\par
``Лера - это не просто человек, который постит. Она знает, когда публиковать, какие хештеги использовать и как Instagram 2026 читает аккаунт. Монтажёр может собрать видео. Но без Леры это видео никто не увидит.''
\end{softbox}
\vspace{2mm}
\begin{tabularx}{\linewidth}{Y Y}
\rowcolor{paperbg}
\textbf{Калиброванный вопрос} & \textbf{NLP-привязка} \\
Что произойдёт с охватом, если хештег-стратегия и тайминг публикаций перестанут быть системными? & Retention в первые секунды reels, saves и DM shares важнее лайков. KPI роста невозможны без системного SMM. \\
\end{tabularx}
\end{glassbox}

\vspace{3mm}

\begin{glassbox}
{\headingfont\fontsize{18}{20}\selectfont Диана}\hfill
{\sansfont\bfseries\color{emerald} Портфолио и статус}

{\color{muted} Главный фотограф}

\vspace{2mm}
\begin{softbox}[colback=paperbg]
\minititle{Функция в системе}\par
Создаёт статусный визуальный актив Fashion Photo - ту самую витрину, которая нужна будущим brand collaborations.
\end{softbox}
\vspace{2mm}
\begin{softbox}[colback=paperbg]
\minititle{Bullet Point для клиента}\par
``Диана - причина, по которой бренды захотят работать с Monica. Fashion Photo - это визитная карточка для Bonpoint, Il Gufo и Baby Dior. Это не то, что можно заменить телефоном. Без Дианы рубрика умирает, а вместе с ней умирает и модельное портфолио для Q3-монетизации.''
\end{softbox}
\vspace{2mm}
\begin{tabularx}{\linewidth}{Y Y}
\rowcolor{paperbg}
\textbf{Калиброванный вопрос} & \textbf{NLP-привязка} \\
Как мы покажем потенциальным брендам профессиональное модельное портфолио без профессионального фотографа? & Fashion Photo раз в 7-14 дней - это статус, визуальная память и актив для медиакита. \\
\end{tabularx}
\end{glassbox}

\clearpage

\begin{glassbox}
{\headingfont\fontsize{22}{24}\selectfont Защита команды - Часть II}\hfill
{\color{muted}\sansfont\fontsize{9.5}{12}\selectfont Сохраняем ритм, ускорение, credibility и каркас системы. Это не ``доп. люди'', это архитектура роста.}
\end{glassbox}

\vspace{3mm}

\begin{glassbox}
{\headingfont\fontsize{18}{20}\selectfont Настя}\hfill
{\sansfont\bfseries\color{emerald} Ежедневный пульс}

{\color{muted} Ведёт страницу Моники}

\vspace{2mm}
\begin{softbox}[colback=paperbg]
\minititle{Функция в системе}\par
Stories, ответы на комментарии и DM, ежедневный engagement с нишей - всё, что даёт сигнал живого аккаунта.
\end{softbox}
\vspace{2mm}
\begin{softbox}[colback=paperbg]
\minititle{Bullet Point для клиента}\par
``Настя - это ежедневный пульс Monica. Если этот ритм остановить хотя бы на 2 недели, алгоритм начнёт забывать аккаунт. Это не расход. Это защита уже сделанной инвестиции.''
\end{softbox}
\vspace{2mm}
\begin{tabularx}{\linewidth}{Y Y}
\rowcolor{paperbg}
\textbf{Калиброванный вопрос} & \textbf{NLP-привязка} \\
Кто будет поддерживать ежедневную алгоритмическую активность аккаунта, если Настя уйдёт? & Пауза в 2-3 недели создаёт ощущение перезапуска. Восстановление потом обходится ощутимо дороже. \\
\end{tabularx}
\end{glassbox}

\vspace{3mm}

\begin{glassbox}
{\headingfont\fontsize{18}{20}\selectfont Настя}\hfill
{\sansfont\bfseries\color{emerald} Ускоритель роста}

{\color{muted} Таргетолог}

\vspace{2mm}
\begin{softbox}[colback=paperbg]
\minititle{Функция в системе}\par
Платный рост, который помогает пройти ``valley of death'' от 279 до 1 000 подписчиков без затяжки на полгода.
\end{softbox}
\vspace{2mm}
\begin{softbox}[colback=paperbg]
\minititle{Bullet Point для клиента}\par
``Настя-таргетолог превращает путь к 1 000 подписчиков из 6 месяцев в 8 недель. Каждая неделя ожидания - это неделя, когда конкуренты занимают внимание твоей аудитории. Она не излишество, она ускоритель.''
\end{softbox}
\vspace{2mm}
\begin{tabularx}{\linewidth}{Y Y}
\rowcolor{paperbg}
\textbf{Калиброванный вопрос} & \textbf{NLP-привязка} \\
Если мы уберём платное продвижение, сколько месяцев ты готова ждать до первых коллабораций? & Без таргета путь к 1K - 6+ месяцев. С таргетом появляется шанс на UGC и первые бренд-запросы уже летом. \\
\end{tabularx}
\end{glassbox}

\vspace{3mm}

\begin{glassbox}
{\headingfont\fontsize{18}{20}\selectfont Маша}\hfill
{\sansfont\bfseries\color{emerald} Credibility и медиа}

{\color{muted} PR}

\vspace{2mm}
\begin{softbox}[colback=paperbg]
\minititle{Функция в системе}\par
Выводит Monica за пределы Instagram: медиа, мероприятия, внешние контакты, брендовая репутация.
\end{softbox}
\vspace{2mm}
\begin{softbox}[colback=paperbg]
\minititle{Bullet Point для клиента}\par
``Маша превращает Instagram-страницу в бренд, о котором говорят. Одно правильное упоминание в медиа или одно верное событие в Монако могут принести не только подписчиков, но и реальные бренды в воронку.''
\end{softbox}
\vspace{2mm}
\begin{tabularx}{\linewidth}{Y Y}
\rowcolor{paperbg}
\textbf{Калиброванный вопрос} & \textbf{NLP-привязка} \\
Как ты планируешь получить первые 2-3 бренд-партнёрства без человека, который отвечает за внешние коммуникации? & PR на ранней стадии отличает проекты, которые взлетают, от тех, что годами стоят на месте. \\
\end{tabularx}
\end{glassbox}

\vspace{3mm}

\begin{glassbox}
{\headingfont\fontsize{18}{20}\selectfont Яна}\hfill
{\sansfont\bfseries\color{emerald} Каркас системы}

{\color{muted} Проджект-менеджер}

\vspace{2mm}
\begin{softbox}[colback=paperbg]
\minititle{Функция в системе}\par
Координация, ритм, защита от хаоса и удержание стратегии как системы, а не набора красивых, но случайных идей.
\end{softbox}
\vspace{2mm}
\begin{softbox}[colback=paperbg]
\minititle{Bullet Point для клиента}\par
``Яна - это причина, почему стратегия существует как система. Без неё 4 рубрики снова станут случайным набором постов, ритм развалится, и проект вернётся к постоянной перестройке.''
\end{softbox}
\vspace{2mm}
\begin{tabularx}{\linewidth}{Y Y}
\rowcolor{paperbg}
\textbf{Калиброванный вопрос} & \textbf{NLP-привязка} \\
Кто будет следить за тем, чтобы 4 рубрики выходили по графику и система не развалилась через 3 недели? & Если убрать PM, клиент экономит на человеке, но снова начинает платить хаосом, переделками и потерей фокуса. \\
\end{tabularx}
\end{glassbox}

\clearpage

\begin{glassbox}
\stephead{7}{Future Pacing}

\vspace{3mm}
\begin{tabularx}{\linewidth}{Y Y}
\rowcolor{paperbg}
\textbf{Сценарий A - Полная команда} & \textbf{Сценарий B - Только трое} \\
\fontsize{10.2}{14}\selectfont
Через 8 недель, к июню 2026, Monica выходит на 800-1000 подписчиков.
Reels стабильно дотягиваются до 5K просмотров.
Bonpoint присылает первый бартер.
У команды уже есть медиакит, который Маша отправила в 10 релевантных редакций.
Диана делает фотосессию, которой хочется гордиться.
Ира доводит страницу до уровня, где бренды сами пишут в DM.
А Яна держит процесс настолько чистым, что клиенту не нужно жить в операционке.
&
\fontsize{10.2}{14}\selectfont
Через 8 недель аккаунт остаётся в районе 350 подписчиков.
Контент выходит нерегулярно, reels без сильного хука набирают 200 просмотров.
Нет PR, нет коллабораций, нет медиакита.
Фото на телефон визуально не отличают Monica от десятков других страниц.
И снова появляется мысль: ``Может, нужна новая стратегия?''
На практике это означает возврат к хаосу и более дорогой перезапуск позже. \\
\end{tabularx}
\end{glassbox}

\vspace{4mm}

\begin{glassbox}
\stephead{8}{Agreement Frame}

\vspace{3mm}
\begin{navybox}
\eyebrow{Финальное предложение}\par
\vspace{2mm}
{\headingfont\fontsize{18}{20}\selectfont 8-недельный спринт вместо резкого сокращения\par}
\vspace{3mm}
``Я понимаю, что бюджет - это важно. И при этом мы обе хотим, чтобы Monica стала брендом, а не просто красивой страницей.
Я предлагаю не сокращать людей сейчас, а зафиксировать 8-недельный спринт.
Через 8 недель мы вместе оценим результат: если рост не соответствует KPI, мы сами предложим оптимизацию.
Но если система покажет результат - а она его покажет - ты будешь рада, что не сократила её слишком рано.
Как тебе такой вариант?''
\end{navybox}

\vspace{3mm}
\begin{multicols}{2}
\begin{softbox}[colback=paperbg]
\textbf{1.} Не говорим ``нет'' сокращению. Говорим: ``Давай оценим честно через 8 недель''.
\end{softbox}
\vspace{2.5mm}
\begin{softbox}[colback=paperbg]
\textbf{2.} Даём контроль через KPI: рост подписчиков, ER, бренд-контакты.
\end{softbox}
\vspace{2.5mm}
\begin{softbox}[colback=paperbg]
\textbf{3.} Используем presupposition: система покажет результат, если её не ломать раньше времени.
\end{softbox}
\vspace{2.5mm}
\begin{softbox}[colback=paperbg]
\textbf{4.} Финальный вопрос мягкий: ``Как тебе такой вариант?'' - клиент соглашается своей формулировкой.
\end{softbox}
\end{multicols}
\end{glassbox}

\vspace{4mm}

\begin{glassbox}
{\headingfont\fontsize{20}{22}\selectfont Сводная таблица}\par
\vspace{2mm}
\rowcolors{2}{paperbg}{paperwhite}
\begin{tabularx}{\linewidth}{>{\raggedright\arraybackslash}p{0.15\linewidth} >{\raggedright\arraybackslash}p{0.18\linewidth} >{\raggedright\arraybackslash}p{0.31\linewidth} Y}
\toprule
\rowcolor{paperbg}
\textbf{Человек} & \textbf{Функция бренда} & \textbf{Что теряем без этого человека} & \textbf{Стоимость потери} \\
\textbf{Ира} & Визуальный язык & Контент становится обычным, исчезает quality overshoot. & Потеря USP и визуального превосходства. \\
\textbf{Лера} & Алгоритм и стратегия & Ломаются хештеги, timing и caption SEO. & Reach может просесть на 50-70\%. \\
\textbf{Настя} & Ежедневный пульс & Алгоритм перестаёт считывать Monica как живой аккаунт. & Восстановление ритма обойдётся дороже текущей работы. \\
\textbf{Диана} & Модельное портфолио & Fashion Photo умирает, бренду нечего показывать партнёрам. & Замедление монетизации и потери collab-воронки. \\
\textbf{Настя (таргет)} & Ускоритель роста & Путь к 1K растягивается на месяцы. & 4+ месяца упущенного импульса. \\
\textbf{Маша} & Credibility & Нет медиа, нет выхода за пределы Instagram. & Годы вместо месяцев до статуса. \\
\textbf{Яна} & Каркас системы & Рубрики распадаются, стратегия снова становится хаосом. & Потеря всей архитектуры роста. \\
\bottomrule
\end{tabularx}
\end{glassbox}

\clearpage

\begin{glassbox}
{\headingfont\fontsize{23}{25}\selectfont Шпаргалка для Яны}\hfill
{\color{muted}\sansfont\fontsize{9.5}{12}\selectfont Краткий боевой лист на переговоры: в каком порядке говорить и какое ощущение создавать у клиента.}
\end{glassbox}

\vspace{3mm}

\begin{multicols}{3}
\begin{softbox}[colback=paperwhite]
\minititle{1. Открытие}\par
Начни с Accusation Audit: перечисли 4 главных страха клиента раньше неё самой.
\end{softbox}

\vspace{2.5mm}
\begin{softbox}[colback=paperwhite]
\minititle{2. Labeling}\par
На каждое возражение отвечай ``Похоже, ты чувствуешь...'' вместо аргумента или спора.
\end{softbox}

\vspace{2.5mm}
\begin{softbox}[colback=paperwhite]
\minititle{3. That's Right}\par
Точно суммируй её позицию и не двигайся дальше, пока не услышишь ``Вот именно''.
\end{softbox}

\vspace{2.5mm}
\begin{softbox}[colback=paperwhite]
\minititle{4. Reframe}\par
Переведи разговор из формулы ``7 человек'' в формулу ``4 функции бренда''.
\end{softbox}

\vspace{2.5mm}
\begin{softbox}[colback=paperwhite]
\minititle{5. Факты}\par
Дай 2-3 сильных сигнала: алгоритм забывает, перезапуск дороже, early-stage бренд нельзя резать раньше времени.
\end{softbox}

\vspace{2.5mm}
\begin{softbox}[colback=paperwhite]
\minititle{6. Команда}\par
Каждую роль объясняй через функцию, боль потери и один калиброванный вопрос.
\end{softbox}

\vspace{2.5mm}
\begin{softbox}[colback=paperwhite]
\minititle{7. Будущее}\par
Нарисуй два сценария через 8 недель: яркий и тусклый. Пускай контраст работает за вас.
\end{softbox}

\vspace{2.5mm}
\begin{softbox}[colback=paperwhite]
\minititle{8. Закрытие}\par
Не упирайся. Предложи 8-недельный спринт с KPI и совместной оценкой по итогам.
\end{softbox}

\vspace{2.5mm}
\begin{softbox}[colback=paperwhite]
\minititle{Тон}\par
Late-Night FM DJ: медленный, тёплый, уверенный. Не давим. Не оправдываемся. Показываем.
\end{softbox}
\end{multicols}

\vspace{4mm}

\begin{glassbox}
{\color{muted}\sansfont\fontsize{9.5}{12}\selectfont
Playbook создан: 1 апреля 2026 \hfill
Источники: Chris Voss (FBI), NLP framework, Reddit, The Drum, CMSwire, MarketingFuel UK, SpurNow, InfluenceFlow 2026
}
\end{glassbox}

\end{document}
